\section{First Fundamental Theorem of Calculus}

Till now Definite and Indefinite integral are defined
very differently as below . 

\begin{tabular}{l|l}
	\toprule
	Definite Integral &           Indefinite Integral\\
	\midrule
	Denoted by :      &                    Denoted by \\
	$\int^{b}_{a} f(x) \; dx$  & $\int f(x) \; dx$  \\

	defined as :      &                 defined as \\

	\import{./images/}{definite_integral.pdf_tex}
						      &	Antiderivative i.e Inverse of differentiation    \\

	Calculated by : 
	Riemann sums      &                            \\
\end{tabular}

\paragraph{Objectives} 

\begin{enumerate}

    \item Statement of FTC1 and use it to evaluate integrals 
    \item  use properties of definite integrals such as additivity , symmetry
    \item use method of substitution to evaluate definite integrals. 

\end{enumerate}

\subsection{ The FTC 1 } 

\begin{definition*}
    The FTC 1 states that , if $F$ is differentiable and $\dot F = f$ is continous , then 
    \begin{equation*}
        \int^{b}_{a} f(x) \; dx = F(b) - F(a) = F(x) \Biggr |_{a}^{b} 
    \end{equation*}
\end{definition*}

The FTC1 connects definite integral to Antiderivative . To be consistent with FTC1 
Riemann sum formula for definite integral doesnot need to be changed . 
Consequently , the geometric definition of $\int^{b}_{a} f(x) \; dx$ that is 
consistent with FTC1 is as follows . 

\begin{figure}[H]
	\centering
	\def\svgwidth{\columnwidth}
	\import{./images/}{FTC.pdf_tex}
\end{figure}

In other words , 
\begin{definition*}
    The definite integral of a function is the signed area bounded by curve 
\end{definition*}

\subsection{Velocity and Speed , Displacement and Distance travelled}

you are travelling between time a and time b with velocity $v(t)$ ,Then 
the displacement be $x(t)$ is given by 
\[
	\int^{b}_{a} v(t) dt = x(b) -x(a)
\]
The total distance travelled is given by 
\[
	\int^{b}_{a} |{v(t)}| \; dx
\]
